\section{What Makes a Compressed Format Computable}
\label{sec:principles}

Most compressed formats cannot be computed on, and the reasons are not accidents
of implementation. A format designed purely to minimize bytes will reach for
whichever coupling buys the most redundancy, and coupling is exactly what a query
engine cannot afford. This section states the five properties that separate the
two cases. They are drawn from building \sys, but none of them is specific to
genomics, and we state them in general terms.

\paragraph{P1. Local expansion.}
Producing the $i$th element must not require having reconstructed the preceding
ones, beyond a bounded running state. This is the property that distinguishes a
representation a fold can walk from a stream that must be replayed from the start.
Long-window LZ77 violates it: a match may reference any byte within a window of
megabytes, so the decoder must keep that history materialized. Adaptive
arithmetic coding violates it more strongly, since the model state at position
$i$ depends on the entire prefix. Grammar compression satisfies it: a rule
expands from a static rulebook, independent of where it appears, and expansion
depth is bounded by the number of grammar rounds rather than by data size. A
fixed-width delta transform satisfies it with a single accumulator.

\paragraph{P2. Side-table-free symbol classification.}
The decoder's inner loop must distinguish a literal from a reference in constant
time, without consulting a data structure. This sounds minor and is not: it is
the difference between a comparison and a hash probe on the hottest instruction
path in the system, executed once per emitted element, tens of billions of times
per whole-genome query. The mechanism is namespace partitioning. If reference IDs
are allocated strictly above every literal value, classification is a range check
on the symbol itself. Formats that assign dictionary IDs from a shared namespace,
or that tag symbols out of band, forfeit this.

\paragraph{P3. Column independence.}
A query must be able to read the fields it needs without paying for the fields it
does not. This is the standard argument for columnar layout in analytical
databases~\cite{abadi2006integrating,boncz2005monetdb}, and it applies with more
force to scientific interchange formats, which are almost universally row-major
text. It is what lets one of our analyses answer a question by reading metadata
columns and decoding \emph{no} traversals at all, and what lets the others skip
the segment DNA payload, which is the largest non-traversal component of the file,
whenever they do not need to spell sequence.

\paragraph{P4. Record independence.}
Records must be individually addressable and independently decodable, which in
practice means an explicit length or offset column alongside the payload. This
buys two things at once. It makes the format partitionable, so $N$ threads can
decode $N$ records with no coordination on the decode path, which is what turns
the fold into a parallel loop. And it bounds memory, because a query can hold one
record's expansion at a time rather than the collection's.

\paragraph{P5. Closure under the domain's symmetries.}
Whatever cheap transformations the domain's queries apply, the compressed
representation should express directly rather than force through an expansion. In
our setting the transformation is reverse complement: reading a segment on the
opposite DNA strand emits its elements in reverse order with every orientation
flipped, and roughly half of all traversal steps in a pangenome are on the reverse
strand. A representation that could not reverse a rule in place would have to
store a mirrored copy of the rulebook, or expand and reverse at query time, and
either would negate the benefit. Encoding a negated reference as ``the reverse
complement of this rule'' makes the symmetry free. The general form of the
property is that the operations a query composes should be closed over the
compressed representation, not only over the expanded one.

\begin{table}[t]
\centering
\small
\setlength{\tabcolsep}{4.2pt}
\caption{Formats against the five properties. \checkmark{} holds, $\circ$ holds
partially or only within a block, \texttimes{} does not hold. The row that matters
is not that our container scores well, but that scoring well is a design choice
with a cost (see text): every property traded away some compression ratio.}
\label{tab:principles}
\begin{tabular}{@{}lccccc@{}}
\toprule
Format & P1 & P2 & P3 & P4 & P5 \\
       & local & class. & col. & rec. & sym. \\
\midrule
GFA text                       & \checkmark & \texttimes & \texttimes & $\circ$ & \texttimes \\
gzip / bgzip~\cite{li2011tabix} & \texttimes & \texttimes & \texttimes & $\circ$ & \texttimes \\
Zstd frame~\cite{collet2021rfc} & \texttimes & \texttimes & \texttimes & $\circ$ & \texttimes \\
Parquet\,+\,Zstd               & $\circ$ & \texttimes & \checkmark & $\circ$ & \texttimes \\
HDF5 / Zarr (chunked)          & $\circ$ & \texttimes & \checkmark & $\circ$ & \texttimes \\
GBZ / GBWT~\cite{siren2022gbz} & \checkmark & \checkmark & \texttimes & \texttimes & \checkmark \\
This container~\cite{gfaz2026} & \checkmark & \checkmark & \checkmark & \checkmark & \checkmark \\
\bottomrule
\end{tabular}
\end{table}

Table~\ref{tab:principles} places common formats against the five properties. Two
readings of it are worth separating. The block-compressed formats in the middle
rows are not disqualified from analytics in general; Parquet and Zarr are built
for it and satisfy P3 well. What they do not offer is P1 and P2 \emph{within} a
block: the decompressor's output is the expanded data, so a query over a Parquet
column materializes that column. That is entirely reasonable when the expanded
column is small relative to memory, and it is the assumption that fails here,
where the expanded traversal column is the $369$\,GB object we are trying to
avoid. GBZ satisfies P1, P2, and P5 and is the closest existing design, but its
index is global rather than per-record, which is what P4 asks for.

\paragraph{Computability is bought, not found.}
The properties are not free, and a paper that presented them as pure wins would be
misleading. Namespace partitioning (P2) reserves a range of the symbol space that
a tighter encoding could have used, and forgoes the shorter codes a fully shared
dictionary would permit. Column independence (P3) forgoes cross-field correlation:
fields that predict one another compress better together, and splitting them
sacrifices that. Record independence (P4) is the most expensive of the five,
because it prevents redundancy from being exploited across record boundaries, and
in a dataset whose entire redundancy structure is cross-record, that is precisely
the redundancy one most wants. The container recovers it by making the grammar
rulebook global while keeping the encoded records independent, so rules are shared
across every traversal but each traversal decodes alone; this is the single
design decision that makes the rest of the engine possible, and we return to it in
Section~\ref{sec:container}.

The upshot is a claim we can state generally. Formats intended for datasets that
will outgrow memory should be designed against these five properties from the
start, accepting a few percent of ratio to keep the data computable, because the
alternative is that every consumer pays full materialization forever.
